\documentclass[letterpaper,10pt]{article}

\usepackage[margin=0.58in, top=0.48in, bottom=0.48in]{geometry}
\usepackage{enumitem}
\usepackage{titlesec}
\usepackage[T1]{fontenc}
\usepackage[hidelinks]{hyperref}
\usepackage{xcolor}

\pagestyle{empty}
\setlength{\parindent}{0pt}
\setlength{\parskip}{0pt}

\titleformat{\section}
  {\normalsize\bfseries}
  {}
  {0em}
  {}
  [\titlerule]

\titlespacing{\section}{0pt}{8pt}{3pt}

\newcommand{\resumeentry}[4]{
  \noindent\textbf{#1} \hfill #2 \\
  \textit{#3} \hfill \textit{#4}
}

\newcommand{\projectentry}[3]{
  \noindent\textbf{#1} \textbar{} \textit{#2} \hfill #3
}

\newcommand{\activityentry}[2]{
  \noindent\textbf{#1} \hfill #2
}

\setlist[itemize]{
  leftmargin=1.25em,
  topsep=1pt,
  itemsep=1pt,
  parsep=0pt,
  partopsep=0pt
}

\begin{document}

\begin{center}
  {\Large\textbf{Amelia Eckard}} \\[3pt]
  Charlotte, NC \textbar{}
  \href{https://ameliaeckard.com}{ameliaeckard.com} \textbar{}
  \href{https://github.com/ameliaeckard}{github.com/ameliaeckard} \textbar{}
  \href{https://linkedin.com/in/ameliaeckard}{linkedin.com/in/ameliaeckard}
\end{center}

\section{Education}

\resumeentry
  {University of North Carolina at Charlotte}{Charlotte, NC}
  {M.S. in Artificial Intelligence \textbar{} Early Entry Program}{Expected May 2028}

\smallskip

\resumeentry
  {University of North Carolina at Charlotte}{Charlotte, NC}
  {B.S. in Computer Science \textbar{} Concentration: AI, Robotics, \& Gaming}{Expected May 2027}

\section{Experience}

\resumeentry
  {Research Intern}{Remote}
  {Stevens Institute of Technology, Dr. Jina Huh-Yoo}{June 2026 -- Present}
\begin{itemize}
  \item Engineered \textit{rq-system}/LEAF, an AI-assisted literature-analysis pipeline using PubMed E-utilities, CrossRef, and LLM APIs to screen, classify, and extract structured evidence from 345 papers; processed 256 included, 80 excluded, and 9 review-flagged papers.
  \item Implemented dataset tagging, a seven-category research-question taxonomy, explicit screening/classification/extraction gates, and gold-sample validation using precision, recall, F1, accuracy, and confusion matrices.
\end{itemize}

\smallskip

\resumeentry
  {Undergraduate Researcher}{Charlotte, NC}
  {UNC Charlotte, Dr. Todd Dobbs}{August 2025 -- May 2026}
\begin{itemize}
  \item Developed an Apple Vision Pro indoor-navigation prototype for visually impaired users using Swift, ARKit, RealityKit, ObjectTrackingProvider, Core ML, HRTF spatial audio, and distance-based pitch cues.
  \item Investigated object tracking, on-device recognition, and spatial-audio guidance as complementary modalities for accessible indoor navigation. Designed evaluation protocols for object identification, tracking reliability, navigation performance, and spatial-audio feedback workflows.
\end{itemize}

\smallskip

\resumeentry
  {Lead Instructional Assistant, ITSC 1213}{Charlotte, NC}
  {UNC Charlotte College of Computing and Informatics}{August 2026 -- Present}
\begin{itemize}
  \item Serve as lead instructional assistant for an asynchronous Python programming course, coordinating grading, rubrics, office hours, student support, and course logistics with the instructor.
\end{itemize}

\section{Projects}

\projectentry
  {CNN Scene Classification}
  {Python, PyTorch, torchvision, CUDA}
  {2026}
\begin{itemize}
  \item Trained and evaluated convolutional neural networks for 16-class scene recognition on 2,400 images.
  \item Improved validation accuracy from 42.25\% to 51.46\% through image augmentation while comparing preprocessing choices, training behavior, and model performance.
\end{itemize}

\smallskip

\projectentry
  {Emergency Fund Predictor}
  {Python, scikit-learn, pandas, SHAP}
  {2025}
\begin{itemize}
  \item Built an 11-feature machine-learning pipeline using 12,295 Federal Reserve SHED respondents enriched with BEA indicators to study emergency-fund vulnerability.
  \item Compared three classifiers and applied SHAP explainability to identify the strongest drivers of financial fragility and model predictions.
\end{itemize}

\smallskip

\projectentry
  {Resolve -- LPL Financial University Hackathon}
  {AI/ML, Workflow Intelligence, FinTech}
  {2026}
\begin{itemize}
  \item Developing a proactive exception-intelligence layer that predicts likely blockers before financial-service requests visibly stall.
  \item Designing case-history and document understanding, workflow-state comparison, anomaly and stall-risk detection, causal explanations, and next-best-action recommendations.
\end{itemize}

\smallskip

\projectentry
  {Vision Lab: ASCII Reconstruction}
  {JavaScript, HTML/CSS, Image Processing}
  {2026 -- Present}
\begin{itemize}
  \item Developed an image-to-ASCII reconstruction experiment for a web-based computer-vision lab connected to ameliaeckard.com.
\end{itemize}

\section{Technical Skills}

\noindent\textbf{Research/AI} \quad PyTorch, scikit-learn, NumPy, pandas, OpenCV, SHAP, Core ML, LLM APIs \\[1pt]
\noindent\textbf{Methods} \quad Precision, recall, F1, accuracy, confusion matrices, structured extraction, taxonomy design \\[1pt]
\noindent\textbf{Tools/APIs} \quad Python, Git, CUDA, PubMed E-utilities, CrossRef, Perplexity API

\section{Certifications \& Involvement}

\activityentry
  {UR2PhD Training Course, Computing Research Association}{December 2025}

\smallskip

\activityentry
  {Social and Behavioral Research, CITI Program}{September 2025}

\end{document}

\smallskip

\activityentry
  {Kode With Klossy \texttimes{} Goldman Sachs Machine Learning Challenge}{December 2025}

\end{document}
